% Part: first-order-logic
% Chapter: models-theories
% Section: expressing-props-of-structures

\documentclass[../../../include/open-logic-section]{subfiles}

\begin{document}

% SEGMENT OLP-0169-S01

\olfileid{fol}{mat}{exs}

\olsection{\printtoken{P}{structure} இன் பண்புகளை வெளிப்படுத்துதல்}

\begin{explain}
சார்புகள் மற்றும் தொடர்புகள் மீதான நிபந்தனைகளை, அல்லது பொதுவாக ஒரு
கட்டமைப்பில் உள்ள சார்புகளும் தொடர்புகளும் இந்நிபந்தனைகளை
நிறைவுறுத்துகின்றன என்பதை வெளிப்படுத்துவது பல நேரங்களில் பயனுள்ளதும்
முக்கியமானதும் ஆகும். எடுத்துக்காட்டாக, !!{predicate}s எதைக்
``குறிக்க'' வேண்டும் என்று நாம் விரும்புகிறோமோ அதை ``பிடிக்கும்'' ஒரு
மொழிக்கான !!{structure}s ஐ, அவ்வாறு செய்யாதவற்றிலிருந்து வேறுபடுத்த
வழிகள் வேண்டும். எந்த !!{structure}s ஐ நாம் ``நோக்கமாகக் கொள்கிறோம்''
என்று குறிப்பிட நமக்கு முழுச் சுதந்திரம் உண்டு. எடுத்துக்காட்டாக,
!!{predicate}~$\le$ இன் பொருள்கோள் ஒரு வரிசையாக இருக்க வேண்டும் என்று
குறிப்பிடலாம். அல்லது பொருட்களம் கணங்களைக் கொண்டிருந்து, $\Obj \in$
ஆனது ``ஓர் !!a{element} ஆக இருத்தல்'' என்ற தொடர்பாகப்
பொருள்கொள்ளப்படும்~$\Lang L$ இன் பொருள்கோள்களில் மட்டுமே ஆர்வம் உள்ளது
என்று குறிப்பிடலாம். ஆனால் மொழியின் !!{sentence}s கொண்டு இதைச் செய்ய
முடியுமா? வேறுவிதமாகச் சொன்னால்,~$\Struct M$ இன் மொழியில்
!!a{sentence} ஒன்றால் (அல்லது !!{sentence}s இன் ஒரு தொகுப்பால்)
!!a{structure}~$\Struct M$ மீதான எந்த நிபந்தனைகளை வெளிப்படுத்தலாம்?
சில நிபந்தனைகளை நாம் வெளிப்படுத்த முடியாது. எடுத்துக்காட்டாக,
$\Domain M = \Nat$ எனில் மட்டுமே ஒரு !!{structure}~$\Struct M$ இல்
மெய்யாகும்~$\Lang L_A$ இன் வாக்கியம் எதுவும் இல்லை.
``பொருட்களம் இயல் எண்களை மட்டுமே கொண்டுள்ளது'' என்பதை வெளிப்படுத்த
முடியாது. ஆனால் !!{structure}s இன் சில ``கட்டமைப்புப் பண்புகளை''
நாம் வெளிப்படுத்தக்கூடும். !!{structure}s இன் எந்தப் பண்புகளை
!!{sentence}s கொண்டு வெளிப்படுத்தலாம்? வேறுவிதமாகச் சொன்னால்,
!!a{sentence} ஒன்றை (அல்லது !!{sentence}s இன் தொகுப்பை) மெய்யாக்கும்
!!{structure}s என எந்தத் தொகுப்புகளை விவரிக்கலாம்?
\end{explain}

\begin{defn}[ஒரு தொகுப்பின் மாதிரி]
$\Gamma$ ஆனது மொழி~$\Lang L$ இன் !!{sentence}s கொண்ட தொகுப்பாக
இருக்கட்டும். எல்லா $!A \in \Gamma$ க்கும் $\Sat{M}{!A}$ எனில்,
!!a{structure}~$\Struct M$ ஆனது $\Gamma$ இன்
\emph{ஒரு மாதிரி} என்கிறோம்.
\end{defn}

\begin{ex}
வாக்கியம் $\lforall[x][x \le x]$ ஆனது~$\Struct M$ இல் மெய் iff
$\Assign{\le}{M}$ ஒரு தற்சுட்டுத் தொடர்பு. வாக்கியம்
$\lforall[x][\lforall[y][((x \le y \land y \le x) \lif x = y)]]$ ஆனது
~$\Struct M$ இல் மெய் iff $\Assign{\le}{M}$ எதிர்சமச்சீர்த் தொடர்பு.
வாக்கியம் $\lforall[x][\lforall[y][\lforall[z][((x \le y \land y \le z)
      \lif x \le z)]]]$ ஆனது~$\Struct M$ இல் மெய் iff $\Assign{\le}{M}$
கடப்புத் தொடர்பு. ஆகவே, பின்வருவதன் மாதிரிகள்
\begin{align*}
\{\quad &\lforall[x][x \le x], \\
   & \lforall[x][\lforall[y][((x \le y \land y \le
    x) \lif x = y)]], \\
   &\lforall[x][\lforall[y][\lforall[z][((x \le y
      \land y \le z) \lif x \le z)]]] \quad \}
\end{align*}
~$\Assign{\le}{M}$ தற்சுட்டு, எதிர்சமச்சீர் மற்றும் கடப்பு ஆகிய
பண்புகளைக் கொண்டிருக்கும் கட்டமைப்புகளே; அதாவது அது ஒரு பகுதி வரிசை.
எனவே இவற்றைப் \emph{பகுதி வரிசைகளின் முதல்தரக் கோட்பாட்டுக்கான}
அடிகோள்களாக எடுத்துக்கொள்ளலாம்.
\end{ex}

\end{document}
